\section{The degree-zero distribution route and its exact remaining gate}
\label{sec:booker-route}

This section examines a further route to the $S_5$ comparison, using
Booker's published distributional converse theorem. It yields a new
necessary signature for the positive-excess branch, but does not
eliminate that branch.

\paragraph{The primary input and its normalization.}
For an $L$-datum $\mathcal F=(f,K,m)$, Booker's normalization is
\[
 L_{\mathcal F}(s)
 =\exp\left(\sum_{n\ge2}\frac{f(n)}{\log n}
                              n^{1/2-s}\right)
 =\sum_{n\ge1}a_{\mathcal F}(n)n^{-s}.
 \tag{MA-LD.1}
\]
Its degree is $2\lim_{x\downarrow0}xK(x)$. Positivity means that
$m(z)<0$ at only finitely many points. It is a condition on the
signed divisor, not on $f(n)$. The exact locators are Definition~1.3,
author p.~3; Examples~1.4(2)--(3), p.~4; Theorem~1.6, p.~5;
and Theorem~1.7, p.~6, in \texttt{arXiv:1308.3067v2}.
The multiplicity-one criteria used below are
\[
 \sum_{n\ge2}\frac{|f(n)|}{\log n}<\infty,
 \qquad
 \sum_{n\ge1}\frac{|a_{\mathcal F}(n)|}{\sqrt n}<\infty,
 \qquad
 \sum_{|\operatorname{Re}z|\le T}|m(z)|=o(T).
 \tag{MA-LD.2}
\]
Each is equivalent to $\mathcal F=0$. In degree zero, the additional
hypothesis of positivity also forces $\mathcal F=0$.
In particular, the second expression in (MA-LD.2) uses the
Dirichlet coefficients $a_{\mathcal F}$, not the logarithmic
coefficients $f$.

\paragraph{MA-LD-001 (an infinite signed discrepancy is necessary).}
Retain all of the family and comparison hypotheses in MA-S5-001,
including $(\mathcal R_B)$. Form the difference of the actual
automorphic and Artin data
\[
 \mathcal F_R=\mathcal F_{\varpi}-\mathcal F_W\in\mathcal L_0,
 \qquad L_{\mathcal F_R}(s)=\frac{A(s)}{L_W(s)}.
 \tag{MA-LD.3}
\]
Let $\mathscr R$ be its completed meromorphic function. It differs
from the quotient $R$ in (MA-S5.3) only by the nonvanishing
conductor exponential and a nonzero normalization constant.
Put
\[
 m_R(z)=\operatorname{ord}_{s=1/2+iz}\mathscr R(s),\qquad
 N_R(T)=\sum_{|\operatorname{Re}z|\le T}|m_R(z)|.
 \tag{MA-LD.4}
\]
If $k>0$, then $\mathscr R$ has infinitely many poles and infinitely
many zeros, all on $\operatorname{Re}s=1/2$, and
\[
                       N_R(T)\ne o(T).
 \tag{MA-LD.5}
\]
The critical-line assertion here concerns the signed quotient
divisor alone and uses no zero-location hypothesis beyond
$(\mathcal R_B)$. It is not a statement about the zeros of the
individual factors.

Consequently, under $(\mathcal R_B)$, each of the following is
equivalent to $k=0$: the quotient has only finitely many poles;
it has only finitely many zeros; or $N_R(T)=o(T)$.

\begin{proof}
Booker's Examples~1.4(2)--(3) attach actual $L$-data to the
cuspidal automorphic representation $\varpi$ and the Artin
representation $W$, respectively. Artin holomorphy is not
needed for the second construction. The class of $L$-data is
an additive group. Both representations have degree five,
so their difference has degree zero. The completed divisor is
the difference of their completed divisors, which proves the
normalizations in (MA-LD.3)--(MA-LD.4). In particular
$m_R(0)=-k$.

For a good prime, logarithmic differentiation of (MA-S5.3)
gives the exact coefficients
\[
 f_R(p^j)=
 \begin{cases}
 4(\log p)p^{-j/2},&p\in\Omega,\ j\text{ odd},\\
 0,&\text{otherwise}.
 \end{cases}
 \tag{MA-LD.6}
\]
The remaining coefficients are supported on powers of finitely
many bad primes. MA-S5-004 shows that
$\pi_\Omega(X)=X^{1/2+o(1)}$. Hence the Euler logarithm in
(MA-LD.1) converges absolutely and locally uniformly on
$\operatorname{Re}s>1/2$: its good first-prime terms are
$4\sum_{p\in\Omega}p^{-s}$, and its higher odd-power terms
are bounded on compact subsets by a constant times
$\sum_p p^{-3\operatorname{Re}s}$.
For each bad prime, axiom~(A1) of Definition~1.3 gives, for
every fixed $a>0$,
$f_R(p^j)(j\log p)^a=O_a(1)$; this also gives the required
local convergence. Exponentiation therefore makes
$L_{\mathcal F_R}$ holomorphic and nonvanishing on that
half-plane.

The divisor $m_R$ is integer-valued everywhere, since it is the
difference of two actual meromorphic divisors. Consequently the
continuation domain in Proposition~2.1(iii) is all of $\mathbb C$;
none of the branch cuts allowed for general $L$-data occurs here.
Proposition~2.1(i), author p.~9, supplies a canonical gamma
factor whose logarithm is holomorphic on
$\operatorname{Re}s>1/2$ and extends continuously to the
boundary. Thus $\mathscr R$ is holomorphic and nonvanishing
on the open right half-plane. Its functional equation,
Proposition~2.1(iv), reflects this conclusion to
$\operatorname{Re}s<1/2$. Its complete signed divisor is
therefore supported on the critical line. Thus their completed
orders cancel off that line. The argument supplies no equality
or inequality between their orders on the critical line.

If $\mathscr R$ had finitely many poles, then
$\mathcal F_R$ would be positive in Booker's precise sense.
The degree-zero case of Theorem~1.7, followed by
Theorem~1.6, would give $\mathcal F_R=0$, contradicting
$m_R(0)=-k<0$. If it had finitely many zeros, apply the same
argument to $-\mathcal F_R$. This proves the two infinitude
assertions. Theorem~1.6(v) gives (MA-LD.5).
Conversely, if $k=0$, MA-S5-001 gives the all-place comparison,
so the difference datum vanishes and all three finiteness
conditions hold.\end{proof}

\paragraph{MA-LD-002 (the summability criteria fail at the exact threshold).}
Under the positive-excess hypotheses of MA-LD-001,
\[
 \sum_{\substack{p\le X\\p\in\Omega}}p^{-1/2}
                  \sim\frac{k}{4}\log\log X,
 \qquad
 \sum_{2\le n\le X}\frac{|f_R(n)|}{\log n}
                  \sim k\log\log X.
 \tag{MA-LD.7}
\]
Also
\[
 \sum_{n\ge1}\frac{|a_{\mathcal F_R}(n)|}{\sqrt n}=\infty,
 \qquad
 \sum_{n\ge2}\frac{|f_R(n)|}{\sqrt n}<\infty.
 \tag{MA-LD.8}
\]
Thus replacing $a_{\mathcal F_R}$ by $f_R$ in Booker's criterion
would produce a false shortcut precisely in this branch.

\begin{proof}
Let
$U(X)=\sum_{p\le X,\ p\in\Omega}(\log p)/\sqrt p$.
MA-S5-004 proves $U(X)\sim(k/4)\log X$. Partial summation gives
\[
 \sum_{\substack{p\le X\\p\in\Omega}}p^{-1/2}
 =\frac{U(X)}{\log X}
       +\int_2^X\frac{U(t)}{t(\log t)^2}\,dt
 \sim\frac{k}{4}\log\log X.
 \tag{MA-LD.9}
\]
The first asymptotic follows by bounding the remainder
$U(t)-(k/4)\log t$ by an arbitrary fixed multiple of
$\log t$ beyond a sufficiently large point. The contribution
of odd powers $j\ge3$ to the second sum in (MA-LD.7) is
bounded by
$4\sum_p\sum_{j\ge3}p^{-j/2}/j<\infty$.
Each bad-prime contribution is finite by (A1), for example
using $|f_R(p^j)|\ll(j\log p)^{-2}$.
The prime term in (MA-LD.6) therefore proves that second
asymptotic.

At each $p\in\Omega$, the Dirichlet coefficient of
$L_{\mathcal F_R}$ is $a_{\mathcal F_R}(p)=4$.
The first assertion of (MA-LD.8) follows already from the
prime subsequence and (MA-LD.7). For the last assertion,
the good-prime contribution is
\[
 4\sum_{p\in\Omega}\sum_{j\text{ odd}}\frac{\log p}{p^j}.
\]
Its higher-power part converges absolutely. For its first
part, $\pi_\Omega(t)\ll t^{3/4}$ eventually, and integration
by parts gives
$\sum_{p\in\Omega}(\log p)/p<\infty$, since
$\int_2^\infty (\log t)t^{-5/4}dt<\infty$.
The finitely many bad-prime geometric tails also converge,
again by (A1).\end{proof}

\paragraph{MA-LD-003 (Euler positivity does not give datum positivity).}
There is an elementary genuine degree-zero example illustrating
the distinction. Let $\chi_5$ be the even primitive quadratic
character modulo $5$, and consider
\[
 L_*(s)=\frac{\zeta(s)}{L(s,\chi_5)},\qquad
 \mathcal F_* =\mathcal F_\zeta-\mathcal F_{\chi_5}.
 \tag{MA-LD.10}
\]
The two archimedean gamma kernels coincide, so $K_*=0$, and
\[
       f_*(n)=(1-\chi_5(n))\frac{\Lambda(n)}{\sqrt n}\ge0
       \quad(n>1).
 \tag{MA-LD.11}
\]
Nevertheless $\mathcal F_*$ is not positive: its completed
quotient has infinitely many poles and infinitely many zeros.

\begin{proof}
The data exist by the usual degree-one examples, and their
difference has degree zero. The logarithmic derivative gives
(MA-LD.11), including the ramified prime $5$, where
$\chi_5(5)=0$. Since $\chi_5(2)=-1$, its coefficient at $2$
is nonzero. If the completed quotient had finitely many
poles, Theorem~1.7 would force this nonzero datum to vanish.
Applying it to the negative datum gives the assertion about
zeros. The quotient has the usual completed functional
equation; the conductor ratio is absorbed in its exponential
factor.\end{proof}

This example is not a number field in the stipulated $S_5$
family and does not have the exceptional-prime sparsity above;
it has a pole at $s=1$. It isolates exactly the invalid general
inference from Euler positivity to positivity of the divisor.

\paragraph{Remaining gate and source-check scope.}
Even holomorphy of both individual factors would not make their
quotient have finitely many poles: a denominator zero creates a
pole whenever its order exceeds that of the numerator. The
degree-zero theorem thus supplies an equivalent cancellation
gate, not a proof of that cancellation. The input $M$ is a
zero-location assertion and supplies no such order inequality.
The ordinary Selberg-class theorem is inapplicable for the same
reason; its analytic-continuation axiom permits a pole only at
$1$, whereas the unresolved branch already has a pole at $1/2$.

The full proofs read for this check were Booker's Lemmas~2.3--2.4,
Proposition~2.1 and its proof, pp.~9--13; the complete
Theorem~1.6 proof in \S2.3, pp.~13--15; and the degree-below-one
part of Theorem~1.7, including Lemma~3.1 and its preceding
contour deformation, pp.~15--17. The exact use of positivity is
on p.~16: it permits a height above all poles before the contour
deformation. Positivity of the coefficients cannot justify that
step. This source check does not claim a new proof of all
automorphic analytic properties used to construct the example
data, nor an audit of the unused degree-one and higher-degree
classification.

Booker, A. R. (2015). $L$-functions as distributions.
\emph{Mathematische Annalen, 363}(1--2), 423--454.
\url{https://doi.org/10.1007/s00208-015-1178-z}.
Author manuscript: \url{https://arxiv.org/abs/1308.3067v2}.
This is a published input, separate from the user-specified
theorem $A$, the conditional claim $M$, and the outside
factoring preprint PF26.
